\documentclass[UTF8, 12pt, fontset=fandol, oneside]{ctexart}
\usepackage{researchnotes}

\graphicspath{{tex/强化学习/figures/math-18/}}

\title{强化学习：前沿}
\author{}
\date{}

\begin{document}

\maketitle

\section*{第 17 章\quad 前沿}

在这最后一章中，我们触及一些超出本书范围、但我们认为对强化学习未来特别重要的主题。
其中许多主题把我们带到可靠已知的内容之外，有些则把我们带到 MDP 框架之外。

\subsection*{17.1\quad 一般价值函数与辅助任务}

贯穿本书，我们关于价值函数的概念已经变得相当一般。
借助离策略学习，我们允许价值函数以任意目标策略为条件。
然后在第 12.8 节中，我们把折扣推广为一个终止函数 $\gamma:\mathcal{S}\mapsto[0,1]$，使得在确定回报（12.17）时，每个时间步都可以应用不同的折扣率。
这使我们能够表达关于在任意的、依赖状态的时域上会得到多少奖励的预测。
下一步，也许是最后一步，是越过奖励，允许对任意信号作预测。
与其预测未来奖励之和，我们可以预测未来声音或颜色感觉值之和，或者预测某种内部的、高度加工过的信号，例如另一个预测。
无论用这种价值函数式预测累加的信号是什么，我们都称之为该预测的累积信号。
我们把它形式化为一个累积信号 $C_t\in\mathbb{R}$。
用它，一个\textbf{一般价值函数}，或 GVF，写作
\begin{equation}
v_{\pi,\gamma,C}(s)=
\mathbb{E}\left[
\sum_{k=t}^{\infty}
\left(\prod_{i=t+1}^{k}\gamma(S_i)\right)C_{k+1}
\mid S_t=s,A_{t:\infty}\sim\pi
\right].
\tag{17.1}
\end{equation}
同常规价值函数（例如 $v_\pi$ 或 $q_*$）一样，这是一个理想函数，我们试图用参数化形式来近似它；我们仍可把近似写作 $\hat v(s,\mathbf{w})$，当然每个预测都必须有不同的 $\mathbf{w}$，也就是说，对 $\pi,\gamma,C$ 的每一种选择都要有一套权重。
因为 GVF 与奖励并没有必然联系，把它称为价值函数也许有些名不副实。
人们也可以直接称它为预测；或者为了让它更有区分度，称为预报。
无论叫什么，它都具有价值函数的形式，因此可以用本书为学习近似价值函数而发展出的通常方法来学习。
与学到的预测一起，我们还可以按通常方式，通过广义策略迭代（第 4.6 节）或演员--评论家方法，学习使这些预测最大化的策略。
这样，一个智能体就可以学习预测并控制大量信号，而不只是长期奖励。

为什么预测和控制长期奖励以外的信号会有用？
这些是\textbf{辅助任务}，因为它们是在最大化奖励这一主任务之外额外存在的任务。
一个回答是，预测和控制许多彼此不同的信号的能力，可以构成一种强大的环境模型。
正如我们在第 8 章中看到的，一个好的模型能使智能体更高效地获得奖励。
要清楚发展这个回答，还需要再引入几个概念，所以我们把它推迟到下一节。
先考虑两种更简单的方式，说明大量多样化预测如何帮助强化学习智能体。

辅助任务帮助主任务的一种简单方式，是它们可能需要主任务同样需要的某些表示。
有些辅助任务可能更容易，延迟更短，动作与结果之间的联系也更清楚。
如果能在容易的辅助任务上早早找到好的特征，那么这些特征就可能显著加快主任务上的学习。
这并没有必然成立的理由，但在许多情形中看起来是合理的。
例如，如果你学会在短时间尺度上，比如几秒内，预测和控制自己的传感器，那么你就可能形成关于物体概念的一部分，而这随后会极大帮助对长期奖励的预测和控制。

我们可以想象一个人工神经网络（ANN），其最后一层被分成多个部分，或称为多个头，每个头负责一个不同任务。
一个头可以产生主任务的近似价值函数（以奖励作为它的累积信号），其余头则产生各种辅助任务的解。
所有头都可以通过随机梯度下降，把误差传播到同一个主体，也就是网络前面共享的部分；这个主体随后会试图在倒数第二层形成能支持所有头的表示。
研究者已经试验过这样的辅助任务，例如预测像素变化、预测下一个时间步的奖励，以及预测回报的分布。
在许多情形中，这种方法已被证明能大幅加快主任务上的学习。
类似地，多重预测也被反复提出，作为引导状态估计构造的一种方式（见第 17.3 节）。

辅助任务学习能改善表现的另一种简单方式，最好通过类比经典条件化这一心理学现象来解释（第 14.2 节）。
理解经典条件化的一种方式是：进化已经把对某个特定信号的预测，与某个特定动作之间建立了一种反射性的（非学习的）关联。
例如，人类和许多其他动物似乎有一个内置反射：只要它们预测眼睛会被戳中的程度超过某个阈值，就眨眼。
预测是学来的，但从预测到闭眼的关联是内置的，因此动物避免了许多眼睛在无保护状态下被戳中的情形。
类似地，从恐惧到心率升高，或到冻结不动的关联，也可以是内置的。
智能体设计者可以做类似的事情：有意地（不通过学习）把关于特定事件的预测连接到预定动作。
例如，一辆自动驾驶汽车如果学会预测继续前行是否会产生碰撞，那么只要这个预测超过某个阈值，就可以给它一个内置反射，让它停车或转向避开。
或者考虑一个吸尘机器人，它学会预测自己在返回充电器之前是否可能耗尽电量；只要这个预测变为非零，它就反射性地驶回充电器。
正确的预测会依赖房子的大小、机器人所在的房间以及电池的老化程度，而这些都很难由机器人设计者预先知道。
设计者很难直接用传感器术语内置一个可靠算法，来决定是否该返回充电器；但如果用学到的预测来表述，这件事可能很容易。
我们预见到许多这样的可能方式：学到的预测可以与内置的行为控制算法有用地结合。

最后，辅助任务也许最重要的作用，是帮助我们超越本书一直采用的一个假设：状态表示是固定的，并且已给定给智能体。
为了解释这个作用，我们首先必须后退几步，体会这一假设的分量，以及去掉它会带来什么影响。
我们在第 17.3 节做这件事。

\subsection*{17.2\quad 通过选项进行时间抽象}

MDP 形式化的一个吸引人之处，是它能有效应用于许多不同时间尺度上的任务。
人们可以用它来形式化这样一些任务：决定收缩哪些肌肉以抓住一个物体，决定乘哪趟航班以方便地到达远方城市，以及决定从事哪份工作以过上满意的人生。
这些任务的时间尺度差异很大，但每一个都可以有用地表述为一个 MDP，并可由本书描述的规划或学习过程来求解。
它们都涉及与世界的交互、序贯决策，以及一个可以有用地理解为随时间累积奖励的目标，因此都可以表述为 MDP。

虽然所有这些任务都可以表述为 MDP，但人们也许会认为它们不能表述为同一个 MDP。
它们涉及如此不同的时间尺度，如此不同的选择和动作概念！
例如，在肌肉抽动的层次上规划横跨大陆的一次飞行是没有用的。
然而对另一些任务，如抓取、投飞镖或击打棒球，低层肌肉抽动也许正是合适的层次。
人们无缝地做所有这些事情，并不显得需要在层次之间切换。
MDP 框架能否被扩展到同时覆盖所有层次？

也许可以。
一个流行思想是：在细致层次上形式化一个 MDP，使用小的时间步，同时允许在更高层次上规划，其中扩展的行动过程对应于许多基础层时间步。
为此，我们需要一种行动过程的概念；它跨越许多时间步，并包含一种终止概念。
把这两个思想形式化的一般方式，是像 GVF 中那样，使用一个策略 $\pi$ 和一个依赖状态的终止函数 $\gamma$。
我们把这样一对对象定义为一种广义动作，称为一个\textbf{选项}。
在时刻 $t$ 执行一个选项 $\omega=(\pi_\omega,\gamma_\omega)$，就是从 $\pi_\omega(\cdot|S_t)$ 得到要采取的动作 $A_t$，然后在时刻 $t+1$ 以概率 $\gamma_\omega(S_{t+1})$ 终止。
如果该选项在 $t+1$ 不终止，那么 $A_{t+1}$ 从 $\pi_\omega(\cdot|S_{t+1})$ 中选出，并且该选项在 $t+2$ 以概率 $\gamma_\omega(S_{t+2})$ 终止，如此继续，直到最终终止。
把低层动作看作选项的特殊情形是方便的：每个动作 $a$ 对应于一个选项 $(\pi_\omega,\gamma_\omega)$，其策略选择该动作（对所有 $s\in\mathcal{S}$，$\pi_\omega(s)=a$），而其终止函数为零（对所有 $s\in\mathcal{S}^+$，$\gamma_\omega(s)=0$）。
选项有效地扩展了动作空间。
智能体可以选择一个低层动作/选项，在一个时间步后终止；也可以选择一个扩展选项，它也许执行许多时间步之后才终止。

选项的设计目标，是让它们可以与低层动作互换。
例如，动作价值函数 $q_\pi$ 的概念自然推广为选项价值函数；它以状态和选项为输入，返回从该状态开始、执行该选项直到终止、随后遵循策略 $\pi$ 所得到的期望回报。
我们还可以把策略的概念推广为\textbf{层级策略}；它从选项而不是动作中选择，其中被选中的选项会一直执行到终止。
有了这些思想，本书中的许多算法都可以推广，以学习近似选项价值函数和层级策略。
在最简单的情形中，学习过程从选项开始“跳到”选项终止，只有在选项终止时才发生更新。
更细致地，也可以在每个时间步都进行更新，使用“选项内”学习算法；这类算法通常需要离策略学习。

选项思想也许使得最重要的推广成为可能：把第 3、4 和 8 章发展的环境模型推广到选项。
一个动作的常规模型，是状态转移概率与每个状态中采取该动作的期望即时奖励。
常规动作模型如何推广成\textbf{选项模型}？
对选项而言，合适的模型同样有两部分：一部分对应于执行该选项所导致的状态转移，另一部分对应于沿途累积奖励的期望。
选项模型中的奖励部分，类似于状态--动作对的期望奖励（3.5），为
\begin{equation}
r(s,\omega)\doteq
\mathbb{E}\left[
R_1+\gamma R_2+\gamma^2R_3+\cdots+\gamma^{\tau-1}R_\tau
\mid S_0=s,A_{0:\tau-1}\sim\pi_\omega,\tau\sim\gamma_\omega
\right],
\tag{17.2}
\end{equation}
其中对所有选项 $\omega$ 和所有状态 $s\in\mathcal{S}$，$\tau$ 是该选项按 $\gamma_\omega$ 终止的随机时间步。
注意这个式子中总体折扣参数 $\gamma$ 的作用：折扣按 $\gamma$ 进行，但选项的终止按 $\gamma_\omega$ 进行。
选项模型的状态转移部分稍微更微妙一些。
模型的这一部分刻画每个可能结果状态的概率（如（3.4）），但现在这个状态可能在各种不同数目的时间步之后才产生，而每一种时间长度都必须以不同方式折扣。
选项 $\omega$ 的模型规定：对 $\omega$ 可能开始执行的每个状态 $s$，以及 $\omega$ 可能终止的每个状态 $s'$，
\begin{equation}
p(s'|s,\omega)\doteq
\sum_{k=1}^{\infty}\gamma^k
\Pr\{S_k=s',\tau=k\mid S_0=s,A_{0:k-1}\sim\pi_\omega,\tau\sim\gamma_\omega\}.
\tag{17.3}
\end{equation}
注意，由于因子 $\gamma^k$ 的存在，这里的 $p(s'|s,\omega)$ 不再是转移概率，也不再在所有 $s'$ 上求和为一。
尽管如此，我们仍继续在 $p$ 中使用“$|$”记号。

上述选项模型转移部分的定义，使我们能够表述适用于所有选项的贝尔曼方程和动态规划算法，其中低层动作是特殊情形。
例如，层级策略 $\pi$ 的状态价值的一般贝尔曼方程为
\begin{equation}
v_\pi(s)=
\sum_{\omega\in\Omega(s)}\pi(\omega|s)
\left[
r(s,\omega)+\sum_{s'}p(s'|s,\omega)v_\pi(s')
\right],
\tag{17.4}
\end{equation}
其中 $\Omega(s)$ 表示状态 $s$ 中可用选项的集合。
如果 $\Omega(s)$ 只包含低层动作，那么这个方程就化为通常贝尔曼方程（3.14）的一个版本；区别只在于新的 $p$（17.3）中包含了 $\gamma$，因此它并不显式出现。
类似地，对应的规划算法也不显式出现 $\gamma$。
例如，与（4.10）类似的选项价值迭代算法为
\[
v_{k+1}(s)\doteq
\max_{\omega\in\Omega(s)}
\left[
r(s,\omega)+\sum_{s'}p(s'|s,\omega)v_k(s')
\right],
\quad \text{对所有 }s\in\mathcal{S}.
\]
如果 $\Omega(s)$ 包含每个 $s$ 中所有可用的低层动作，那么该算法收敛到常规的 $v_*$，由此可以计算最优策略。
然而，当每个状态中只考虑可能选项的一个子集（在 $\Omega(s)$ 中）时，用选项进行规划特别有用。
于是，价值迭代会收敛到受限选项集合内最好的层级策略。
尽管这个策略可能是次优的，但收敛可以快得多，因为考虑的选项更少，也因为每个选项都可以跨过许多时间步。

要用选项进行规划，必须已经给定选项模型，或者学习这些模型。
学习选项模型的一种自然方式，是把它表述为一组 GVF（如前一节定义），然后用本书给出的方法学习这些 GVF。
不难看出，选项模型的奖励部分可以怎样做到这一点。
只需选择一个 GVF，使它的累积信号为奖励（$C_t=R_t$），它的策略为选项的策略（$\pi=\pi_\omega$），而它的终止函数为折扣率乘以选项的终止函数（$\gamma(s)=\gamma\cdot\gamma_\omega(s)$）。
那么真实 GVF 就等于选项模型的奖励部分，$v_{\pi,\gamma,C}(s)=r(s,\omega)$，而本书描述的学习方法可用来近似它。
选项模型的状态转移部分只稍微复杂一些。
需要为该选项可能终止的每一个状态分配一个 GVF。
我们不希望这些 GVF 在选项终止之外累积任何东西，并且只有在终止处于合适状态时才累积。
这可以通过选择预测转移到状态 $s'$ 的 GVF 的累积信号为 $C_t=\gamma(S_t)\cdot\mathbf{1}_{S_t=s'}$ 来实现。
该 GVF 的策略和终止函数选择得与选项模型奖励部分相同。
那么真实 GVF 就等于选项状态转移模型的 $s'$ 部分，$v_{\pi,\gamma,C}(s)=p(s'|s,\omega)$；本书的方法同样可以用来学习它。
虽然每一步看起来都很自然，把它们全部组合起来（包括函数近似和其他必要组成部分）仍相当有挑战，并且超出了当前技术水平。

\noindent\textbf{练习 17.1}\quad 本节为折扣情形给出了选项，但在使用函数近似进行控制时，折扣可以说是不合适的（第 10.4 节）。
对于平均奖励设定（第 10.3 节），与（17.4）类似的层级策略的自然贝尔曼方程是什么？
对于平均奖励设定，与（17.2）和（17.3）类似的选项模型的两个部分是什么？

\subsection*{17.3\quad 观测与状态}

贯穿本书，我们把学到的近似价值函数（以及第 13 章中的策略）写成环境状态的函数。
这是第一部分中方法的一个显著限制；在那里，学到的价值函数实现为一张表，使得任何价值函数都可以被精确近似。
那种情形等价于假设环境状态被智能体完全观测到。
但是，在许多我们感兴趣的情形中，当然也在所有自然智能体的生活中，感觉输入只给出关于世界状态的部分信息。
有些物体可能被其他物体遮挡，或在智能体身后，或在数英里之外。
在这些情形中，环境状态中潜在重要的方面并不能被直接观测到；假设学到的价值函数实现为环境状态空间上的表，是一个强的、不现实的、并且限制性的假设。

我们在第二部分中发展的参数化函数近似框架限制小得多，甚至可以说根本不是限制。
在第二部分中，我们仍保留学到的价值函数（以及策略）是环境状态函数的假设，但允许这些函数由参数化任意地加以限制。
令人稍感意外、且并未被广泛认识到的是，函数近似包含了部分可观测性的重要方面。
例如，如果有一个状态变量不可观测，那么可以选择参数化，使近似价值不依赖该状态变量。
其效果就像这个状态变量不可观测一样。
因此，对参数化情形得到的所有结果，都可以原封不动地应用于部分可观测性。
在这个意义上，参数化函数近似的情形包含了部分可观测性的情形。

不过，还有许多问题如果不更显式地处理部分可观测性，就无法研究。
虽然我们不能在这里完整处理它们，但可以概述需要做出的改变。
一共有四步。

第一，我们会改变问题。
环境不发出它的状态，而只发出\textbf{观测}，也就是依赖其状态、但像机器人的传感器那样只给出部分信息的信号。
为了方便，并且不失一般性，假设奖励是观测的一个直接的、已知的函数（也许观测是一个向量，而奖励是其中一个分量）。
于是，环境交互不再有显式的状态或奖励，而只是动作 $A_t\in\mathcal{A}$ 与观测 $O_t\in\mathcal{O}$ 交替出现的序列：
\[
A_0,O_1,A_1,O_2,A_2,O_3,A_3,O_4,\ldots,
\]
它可以永远继续下去（参见方程 3.1），也可以形成若干幕，每一幕都以一个特殊的终止观测结束。

第二，我们可以从观测和动作序列中恢复本书所使用的状态思想。
我们用词语“历史”和记号 $H_t$ 表示直到某个观测为止的轨迹初始部分：
\[
H_t\doteq A_0,O_1,\ldots,A_{t-1},O_t.
\]
历史表示在不查看数据流之外任何内容的情况下，我们关于过去所能知道的最多信息（因为历史就是整个过去的数据流）。
当然，历史会随着 $t$ 增长而变得巨大且难以处理。
状态的思想是历史的某种紧凑摘要，并且它在预测未来方面与真实历史同样有用。
我们要清楚说明这究竟意味着什么。
要成为历史的摘要，状态必须是历史的函数，$S_t=f(H_t)$；要像完整历史一样有利于预测未来，它必须具有所谓的\textbf{马尔可夫性质}。
形式上，这是函数 $f$ 的一个性质。
当且仅当任意两个历史 $h$ 和 $h'$ 被 $f$ 映射到同一个状态（$f(h)=f(h')$）时，它们对下一观测也具有相同概率，一个函数 $f$ 才具有马尔可夫性质：
\begin{equation}
f(h)=f(h')\Rightarrow
\Pr\{O_{t+1}=o|H_t=h,A_t=a\}
=
\Pr\{O_{t+1}=o|H_t=h',A_t=a\},
\tag{17.5}
\end{equation}
对所有 $o\in\mathcal{O}$ 和 $a\in\mathcal{A}$ 成立。
如果 $f$ 是马尔可夫的，那么 $S_t=f(H_t)$ 就是本书中所用意义下的一个状态。
从现在起，为了把它同那些虽是历史摘要但达不到马尔可夫性质的状态区分开来，我们称它为\textbf{马尔可夫状态}（后者很快就会讨论）。

马尔可夫状态是预测下一观测（17.5）的良好基础；但更重要的是，它也是预测或控制任何事物的良好基础。
例如，令一个\textbf{测试}为未来可能发生的任意一个动作与观测交替出现的特定序列。
例如，一个三步测试记为 $\tau=a_1o_1a_2,o_2,a_3,o_3$。
给定特定历史 $h$ 时，该测试的概率定义为
\begin{equation}
p(\tau|h)\doteq
\Pr\{O_{t+1}=o_1,O_{t+2}=o_2,O_{t+3}=o_3
\mid H_t=h,A_t=a_1,A_{t+1}=a_2,A_{t+2}=a_3\}.
\tag{17.6}
\end{equation}
如果 $f$ 是马尔可夫的，而 $h$ 和 $h'$ 是任意两个在 $f$ 下映射到同一状态的历史，那么任意长度的任意测试 $\tau$，在这两个历史下的概率也必须相同：
\begin{equation}
f(h)=f(h')\Rightarrow p(\tau|h)=p(\tau|h').
\tag{17.7}
\end{equation}
换句话说，马尔可夫状态概括了历史中为确定任何测试概率所需的全部信息。
事实上，它概括了进行任何预测所需的全部信息，包括任何 GVF；也概括了最优行动所需的全部信息（如果 $f$ 是马尔可夫的，那么总有一个确定性函数 $\pi$，使得选择 $A_t=\pi(f(H_t))$ 是最优的）。

把强化学习扩展到部分可观测性的第三步，是处理某些计算方面的考虑。
特别地，我们希望状态是历史的紧凑摘要。
例如，恒等函数完全满足马尔可夫 $f$ 的条件，但仍然用处很小；因为相应状态 $S_t=H_t$ 会随时间增长而变得难以处理，如前所述，更根本的是因为它永远不会重现；在持续任务中，智能体永远不会两次遇到同一个状态，因此也永远无法从表格型学习方法中受益。
我们希望状态既紧凑又马尔可夫。
关于状态如何得到和更新，也有类似问题。
我们并不真正想要一个接受完整历史的函数 $f$。
相反，出于计算原因，我们更愿意通过一种增量、递归的更新获得与 $f$ 相同的效果：它从 $S_t$ 计算 $S_{t+1}$，并纳入下一份增量数据 $A_t$ 和 $O_{t+1}$：
\begin{equation}
S_{t+1}\doteq u(S_t,A_t,O_{t+1}),\quad \text{对所有 }t\ge 0,
\tag{17.8}
\end{equation}
其中给定第一个状态 $S_0$。
函数 $u$ 称为\textbf{状态更新函数}。
例如，如果 $f$ 是恒等函数（$S_t=H_t$），那么 $u$ 只需把 $A_t$ 和 $O_{t+1}$ 追加到 $S_t$ 后面。
给定 $f$ 时，总可以构造一个对应的 $u$，但它在计算上可能并不方便；而且如恒等例子所示，它也可能并不产生紧凑状态。
状态更新函数是任何处理部分可观测性的智能体架构的核心部分。
它必须能高效计算，因为在状态可用之前，不能采取动作，也不能作预测。
这样一个智能体架构的总图如图 17.1 所示。

通过状态更新函数获得马尔可夫状态的一个例子，是流行的贝叶斯方法，称为\textbf{部分可观测 MDP}，或 POMDP。
在这种方法中，环境被假设具有一个定义良好的\textbf{潜在状态} $X_t$；它位于环境观测之下并产生这些观测，但智能体永远得不到它（也不要把它同智能体用于预测和决策的状态 $S_t$ 混淆）。
对 POMDP 而言，自然的马尔可夫状态 $S_t$ 是在给定历史时潜在状态上的分布，称为\textbf{信念状态}。
具体起见，假设通常情形中隐藏状态数有限，$X_t\in\{1,2,\ldots,d\}$。
那么信念状态就是向量 $S_t=\mathbf{s}_t\in\mathbb{R}^d$，其分量为
\[
\mathbf{s}_t[i]\doteq \Pr\{X_t=i\mid H_t\},
\quad \text{对所有可能潜在状态 }i\in\{1,2,\ldots,d\}.
\]
无论 $t$ 如何增长，信念状态都保持同样大小（同样数量的分量）。
假设我们完全知道环境内部如何运作，它还可以由贝叶斯规则增量更新。
具体地说，信念状态更新函数的第 $i$ 个分量为
\begin{equation}
u(\mathbf{s},a,o)[i]\doteq
\frac{\sum_{x=1}^{d}\mathbf{s}[x]p(i,o|x,a)}
{\sum_{x=1}^{d}\sum_{x'=1}^{d}\mathbf{s}[x]p(x',o|x,a)},
\tag{17.9}
\end{equation}
对所有 $a\in\mathcal{A}$、$o\in\mathcal{O}$ 以及分量为 $\mathbf{s}[x]$ 的信念状态 $\mathbf{s}\in\mathbb{R}^d$ 成立；这里的四参数函数 $p$ 不是 MDP 中通常的那个（如第 3 章），而是 POMDP 中用潜在状态表示的类似函数：
\[
p(x',o|x,a)\doteq \Pr\{X_t=x',O_t=o\mid X_{t-1}=x,A_{t-1}=a\}.
\]
这种方法在理论工作中很流行，并且有许多重要应用；但它的假设与计算复杂度伸缩性很差，因此我们不推荐把它作为人工智能的方法。

马尔可夫状态的另一个例子由\textbf{预测状态表示}，或 PSR，给出。
PSR 处理了 POMDP 方法的一个弱点：POMDP 中智能体状态 $S_t$ 的语义扎根在环境状态 $X_t$ 中，而 $X_t$ 从不被观测，因此很难学习。
在 PSR 及相关方法中，智能体状态的语义转而扎根于关于未来观测和动作的预测，而这些量是容易观测的。
在 PSR 中，马尔可夫状态定义为一个 $d$ 维向量，其分量是上面（17.6）定义的 $d$ 个特别选择的“核心”测试的概率。
然后，这个向量由一个状态更新函数 $u$ 更新；该函数类似于贝叶斯规则，但语义扎根于可观测数据，可以说更容易学习。
这种方法已经以许多方式扩展，包括端测试、组合测试、强大的“谱”方法，以及由 TD 方法学习的闭环和时间抽象测试。
一些最好的理论发展是针对称为\textbf{可观测算子模型}（OOM）和序贯系统的系统给出的。

\begin{figure}[htbp]
\centering
\includegraphics[width=0.58\textwidth]{figure-17-1.png}
\caption*{\textbf{图 17.1：}一个概念性智能体架构，包括模型、规划器和状态更新函数。世界在这里接收动作 $A$ 并发出观测 $O$。观测和状态更新函数 $u$ 使用的动作副本产生新状态。新状态输入策略和价值函数，产生下一个动作，同时也输入规划器（以及 $u$）。最负责学习的信息流用虚线表示，这些虚线斜穿过它们所改变的方框。奖励 $R$ 直接改变策略和价值函数。动作、奖励和状态改变模型；模型与规划器紧密合作，也改变策略和价值函数。注意，规划器的运作可以同智能体--环境交互解耦；而其他过程应当与这种交互同步运行，以跟上到来的新数据。还要注意，模型和规划器并不直接处理观测，而只处理由 $u$ 产生的状态；这些状态可以作为模型学习的目标。}
\end{figure}

我们关于如何在强化学习中处理部分可观测性的简要纲要中，第四步也是最后一步，是重新引入近似。
正如第二部分导言所讨论的，要雄心勃勃地接近人工智能，就必须拥抱近似。
这对状态来说同对价值函数来说一样真实。
我们必须接受并使用近似的状态概念。
近似状态将在算法中扮演与以前相同的角色，因此我们继续用记号 $S_t$ 表示智能体所使用的状态，尽管它也许并不是马尔可夫的。

也许最简单的近似状态例子，就是最近的观测，$S_t\doteq O_t$。
当然，这种方法不能处理任何隐藏状态信息。
更好的做法是使用最近 $k$ 个观测和动作：
\[
S_t\doteq O_t,A_{t-1},O_{t-1},\ldots,A_{t-k},
\]
其中某个 $k\ge 1$；这可以通过一个状态更新函数实现，该函数只是把新数据移入、把最旧数据移出。
这种\textbf{$k$ 阶历史}方法仍然很简单，但与试图直接用单个即时观测作为状态相比，它能大幅提高智能体的能力。

当马尔可夫性质（17.5）只近似满足时会发生什么？
不幸的是，当定义马尔可夫性质的一步预测哪怕只是稍有不准确时，长期预测表现也可能急剧恶化。
更长期的测试、GVF 和状态更新函数都可能近似得很差。
短期与长期的近似目标就是不同的，目前还没有有用的理论保证。

尽管如此，仍有理由认为本节概述的一般思想适用于近似情形。
一般思想是：对某些预测有用的状态，对其他预测也有用（特别地，对一步预测充分的马尔可夫状态，也对所有其他预测充分）。
如果我们从马尔可夫情形的这个具体结果后退一步，那么一般思想就类似于第 17.1 节中关于多头学习和辅助任务的讨论。
我们讨论过，适合辅助任务的表示常常也适合主任务。
合起来看，这些观点提出了一种同时处理部分可观测性和表示学习的方法：追求多个预测，并用它们来指导状态特征的构造。
完美但不可行的马尔可夫性质所提供的保证，被一个启发式想法取代：对某些预测有用的东西，也许对其他预测有用。
这种方法能很好地随计算资源伸缩。
有了大型机器，人们可以试验大量预测；也许偏向那些与最终关心的预测最相似的预测，或那些最容易可靠学习的预测，或按某些其他标准来选择。
这里重要的是超越人工选择预测。
智能体应该自己做这件事。
这将需要一种预测的通用语言，使智能体能够系统地探索一个巨大的可能预测空间，筛选出其中最有用的预测。

特别地，POMDP 和 PSR 方法都可以与近似状态一起使用。
状态的语义有助于形成状态更新函数，在这两种方法和 $k$ 阶方法中都是如此。
为了在状态中保留有用信息，语义并不强烈需要是正确的。
有些状态增强方法，例如回声状态网络，会保留关于历史的几乎任意信息，却仍然能够表现良好。
这里有许多可能性，我们期待这一领域出现更多工作和想法。
学习近似状态的状态更新函数，是强化学习中表示学习问题的一个主要部分。

\subsection*{17.4\quad 设计奖励信号}

强化学习相对于监督学习的一个主要优势，是强化学习不依赖详细的教学信息：生成奖励信号并不依赖于知道智能体的正确动作应该是什么。
但是，强化学习应用的成功强烈依赖于奖励信号能多好地框定应用设计者的目标，以及这个信号能多好地评估达到该目标的进展。
由于这些原因，设计奖励信号是任何强化学习应用中的关键部分。

所谓设计奖励信号，我们指的是设计智能体环境中负责在每个时刻 $t$ 计算每个标量奖励 $R_t$ 并把它发送给智能体的那一部分。
在第 14 章末尾关于术语的讨论中，我们说过，$R_t$ 更像是动物大脑内部生成的信号，而不是动物外部环境中的物体或事件。
我们大脑中为我们生成这些信号的部分，经过了数百万年的演化，才很好地适应了祖先为把基因传播给后代而不得不面对的挑战。
因此，我们不应该认为设计一个好的奖励信号总是一件容易的事！

一个挑战是设计奖励信号，使得随着智能体学习，它的行为会接近，并且理想情况下最终达到，应用设计者真正想要的东西。
如果设计者的目标简单且容易识别，例如找到一个定义良好问题的解，或者在定义良好的游戏中获得高分，那么这可以很容易。
在这类情形中，通常会根据智能体成功解决问题的程度，或根据它提高分数的程度来奖励智能体。
但是，有些问题的目标很难转化为奖励信号。
当问题要求智能体熟练执行一项复杂任务或一组任务时，这尤其真实；例如，一个有用的家用机器人助手就需要这样的能力。
进一步说，强化学习智能体可能发现一些出乎意料的方式，让环境给出奖励；其中一些方式也许并不合意，甚至危险。
对任何像强化学习这样基于优化的方法而言，这都是一个长期存在且关键的挑战。
我们会在第 17.6 节，也就是本书最后一节，更多讨论这个问题。

即使存在一个简单且容易识别的目标，也常常会出现\textbf{稀疏奖励}问题。
要足够频繁地给出非零奖励，使智能体至少能达到目标一次，更不用说从多个初始条件高效地学会达到目标，可能是一个令人畏惧的挑战。
显然值得触发奖励的状态--动作对可能相距很远；而标记朝向目标进展的奖励也可能不频繁，因为进展很难或甚至不可能检测。
智能体可能长时间漫无目的地徘徊（Minsky, 1961 称之为“高原问题”）。

在实践中，设计奖励信号常常被留给一种非正式的试错搜索：寻找能产生可接受结果的信号。
如果智能体学不会、学得太慢，或学到错误东西，那么设计者就调整奖励信号并再次尝试。
为此，设计者会用他或她正试图转化为奖励信号的标准来判断智能体表现，使智能体目标同自己的目标相匹配。
如果学习太慢，设计者也许会试图设计一个非稀疏奖励信号，使其在智能体与环境交互的全过程中有效引导学习。

一种诱人的做法，是通过奖励智能体达到设计者认为通向总体目标的重要子目标，来处理稀疏奖励问题。
但是，用善意的补充奖励增强奖励信号，可能导致智能体以非常不同于预期的方式行动；智能体最终可能根本没有达到总体目标。
提供这类引导的更好方式，是保持奖励信号本身不变，而用对价值函数近似最终应该是什么的初始猜测来增强价值函数近似，或者用对它某些部分应该是什么的初始猜测来增强它。
例如，假设想给出 $v_0:\mathcal{S}\to\mathbb{R}$，作为真实最优价值函数 $v_*$ 的初始猜测，并且使用带特征 $\mathbf{x}:\mathcal{S}\to\mathbb{R}^d$ 的线性函数近似。
那么可以把初始价值函数近似定义为
\begin{equation}
\hat v(s,\mathbf{w})\doteq \mathbf{w}^\top\mathbf{x}(s)+v_0(s),
\tag{17.10}
\end{equation}
并像往常一样更新权重 $\mathbf{w}$。
如果初始权重向量为 $0$，那么初始价值函数就是 $v_0$；但渐近解的质量仍像往常一样由特征向量决定。
这种初始化可以用于任意非线性近似器和任意形式的 $v_0$，不过并不能保证总是加快学习。

处理稀疏奖励问题的一种特别有效的方法，是心理学家 B. F. Skinner 引入、并在第 14.3 节描述的\textbf{塑形}技术。
这项技术的有效性依赖于这样一个事实：稀疏奖励问题并不只是奖励信号的问题；它们也是智能体策略的问题，即策略使智能体无法频繁遇到有奖励的状态。
塑形是在学习推进过程中改变奖励信号：从一个在智能体初始行为下并不稀疏的奖励信号开始，逐渐把它修改为适合原本感兴趣问题的奖励信号。
每一次修改都让智能体在其当前行为下仍能频繁得到奖励。
智能体面对的是一串越来越难的强化学习问题；每个阶段学到的东西使下一个更难的问题相对容易，因为相比没有更容易问题的先前经验，智能体现在会更频繁地遇到奖励。
这种塑形是动物训练中的一种基本技术，在计算强化学习中也同样有效。

如果我们完全不知道奖励应该是什么，但有另一个智能体，也许是一个人，已经是该任务的专家，并且其行为可以被观察到，那该怎么办？
在这种情形中，可以使用各种称为“模仿学习”“从示范学习”和“学徒学习”的方法。
这里的思想是从专家智能体中获益，但保留最终做得更好的可能性。
从专家行为学习，可以通过监督学习直接学习，也可以通过所谓\textbf{逆强化学习}抽取奖励信号，然后用强化学习算法配合该奖励信号来学习策略。
Ng 和 Russell（2000）所探索的逆强化学习任务，是试图仅从专家行为中恢复专家的奖励信号。
这无法精确完成，因为一个策略可以相对于许多不同奖励信号都是最优的（例如，任何对所有状态和动作给出同一奖励的奖励信号）；但可以找到合理的候选奖励信号。
不幸的是，这需要很强的假设，包括知道环境动态，以及奖励信号关于哪些特征向量是线性的。
该方法还需要多次完全求解问题（例如用动态规划方法）。
尽管有这些困难，Abbeel 和 Ng（2004）认为，为了从专家行为中获益，逆强化学习方法有时可以比监督学习更有效。

寻找好奖励信号的另一种方法，是把上面提到的对好信号的试错搜索自动化。
从应用角度看，奖励信号是学习算法的一个参数。
同其他算法参数一样，可以通过定义一个可行候选空间并应用优化算法，来自动搜索好的奖励信号。
优化算法用某个候选奖励信号运行强化学习系统若干步，然后用一个旨在编码设计者真实目标的“高层”目标函数给总体结果打分，而忽略智能体的局限。
奖励信号甚至可以通过在线梯度上升来改进，其中梯度是高层目标函数的梯度。
把这种方法同自然界联系起来，优化高层目标函数的算法类似于演化；而高层目标函数则是动物的演化适应度，由其能存活到繁殖年龄的后代数量决定。

使用这种双层优化方法的计算实验已经证实，仅靠直觉并不总是足以设计出好奖励信号；其中一层类似于演化，另一层则归因于个体智能体的强化学习。
当按高层目标函数评估强化学习智能体表现时，结果会以微妙方式对智能体奖励信号的细节非常敏感，而这些方式由智能体的局限以及它行动和学习的环境决定。
这些实验还表明，智能体的目标并不总应与智能体设计者的目标相同。

起初这似乎违反直觉，但无论奖励信号是什么，智能体都可能无法达到设计者的目标。
智能体必须在各种约束下学习，例如计算能力有限、获得关于环境的信息有限，或学习时间有限。
当存在这类约束时，学会达到一个不同于设计者目标的目标，有时最终会比直接追求设计者目标更接近设计者目标。
在自然界中很容易找到这种例子。
由于我们无法直接评估大多数食物的营养价值，演化，即我们奖励信号的设计者，给了我们一个让我们寻求某些味道的奖励信号。
虽然这当然并非万无一失（事实上，在某些与祖先环境有特定差异的环境中可能有害），但它补偿了我们的许多局限：有限的感觉能力、有限的学习时间，以及通过个人实验寻找健康饮食所涉及的风险。
类似地，因为动物不能观察自己的演化适应度，所以这个目标函数不能作为学习的奖励信号。
演化转而提供对演化适应度的可观测预测因子敏感的奖励信号。

最后，请记住，强化学习智能体不一定像一个完整的有机体或机器人；它可以是一个更大的行为系统中的组成部分。
这意味着奖励信号可能受到更大行为智能体内部事物的影响，例如动机状态、记忆、想法，甚至幻觉。
奖励信号也可能依赖学习过程本身的性质，例如衡量学习取得了多少进展的量。
让奖励信号对这类内部因素的信息敏感，使智能体能够学会如何控制它所属的“认知架构”，并获得一些如果奖励信号只依赖外部事件就难以学到的知识和技能。
这些可能性导向了“内在动机强化学习”的思想，我们会在下一节末尾进一步简要讨论。

\subsection*{17.5\quad 剩余问题}

在本书中，我们介绍了强化学习通向人工智能的一种方法的基础。
粗略地说，这种方法基于免模型方法和基于模型方法的共同工作，如第 8 章的 Dyna 架构；并结合第二部分发展的函数近似。
重点一直放在线上、增量算法上；我们认为即使对基于模型的方法来说，这也是根本性的；同时也关注这些算法如何应用于离策略训练情形。
后一件事的完整理由只在最后这一章中给出。
也就是说，我们一直把离策略学习作为处理探索/利用困境的一种有吸引力的方式来介绍；但只有在本章中，我们才讨论了用 GVF 同时学习许多不同的辅助任务，以及按时间抽象的选项模型以层级方式学习世界，这二者都涉及离策略学习。
仍有许多工作要做；正如我们贯穿全书所指出的，也正如本章讨论的额外研究方向所表明的那样。
但是，假设我们宽容一些，承认本书已经完成的一切以及本章到目前为止概述的一切的大致轮廓。
那之后还剩下什么？
当然，我们无法确切知道未来需要什么，但可以作一些猜测。
本节强调六个我们看来仍需由未来研究处理的进一步问题。

第一，我们仍然需要强大的参数化函数近似方法，它们能在完全增量和线上设定中良好工作。
基于深度学习和人工神经网络的方法朝这个方向迈出了一大步；但它们仍然只在大数据集上的批量训练、大规模离线自我博弈训练，或从同一任务上多个智能体交错经验中学习时表现良好。
这些以及其他设定，都是绕开当今深度学习方法一个基本限制的办法：这些方法很难在本书强调的、对强化学习算法最自然的增量线上设定中快速学习。
这个问题有时被描述为“灾难性干扰”或“相关数据”问题。
当学到新东西时，它往往会替代先前学到的东西，而不是添加到其上，结果是较早学习的好处丢失。
诸如“回放缓冲区”这样的技术常被用来保留并重放旧数据，使其好处不至于永久丢失。
诚实的评价必须是：当前深度学习方法并不很适合线上学习。
我们看不出这个限制不可克服的理由，但既能解决它、又保留深度学习优势的算法尚未被设计出来。
当前大多数深度学习研究都在绕开这个限制，而不是消除它。

第二（也许与第一点密切相关），我们仍然需要学习特征的方法，使后续学习能够很好泛化。
这个问题是一个一般问题的实例，它有各种名称，如“表示学习”“构造性归纳”和“元学习”：怎样利用经验，不只是学习一个给定的期望函数，而是学习归纳偏置，使未来学习泛化得更好并因此更快？
这是一个老问题，可以追溯到 20 世纪 50 和 60 年代人工智能与模式识别的起源。\footnote{有人会声称深度学习解决了这个问题，例如第 16.5 节描述的 DQN 展示了一个解，但我们并不信服。目前几乎没有证据表明，仅靠深度学习就以一般且高效的方式解决了表示学习问题。}
这样久远的问题应当让人谨慎。
也许它没有解。
但同样可能的是，找到解并展示其有效性的时机尚未到来。
今天机器学习的规模远大于过去，一个好的表示学习方法的潜在收益已经变得更加明显。
我们注意到，一个新的年度会议，即国际学习表示会议，自 2013 年以来每年都在探索这个及相关主题。
在强化学习语境中探索表示学习则不那么常见。
强化学习为这个老问题带来了一些新的可能性，例如第 17.1 节讨论的辅助任务。
在强化学习中，表示学习问题可以等同于学习第 17.3 节讨论的状态更新函数的问题。

第三，我们仍然需要可伸缩的方法，用学到的环境模型进行规划。
在 AlphaGo Zero 和计算机国际象棋等应用中，规划方法已经被证明极其有效；在那里，环境模型来自游戏规则，或可以由人类设计者以其他方式提供。
但是，完整的基于模型的强化学习，即环境模型从数据中学习然后用于规划的情形，仍然罕见。
第 8 章描述的 Dyna 系统是一个例子；但正如那里以及大多数后续工作所描述的，它使用的是没有函数近似的表格型模型，这极大限制了它的适用性。
只有少数研究包含学到的线性模型，更少的研究还探索了包含使用选项的时间抽象模型，如第 17.2 节所讨论的那样。

在用学到模型进行规划能够有效之前，还需要更多工作。
例如，模型学习需要具有选择性，因为模型的范围会强烈影响规划效率。
如果一个模型聚焦于最重要选项的关键后果，那么规划可以高效且快速；但如果模型包含那些不太可能被选择的选项的不重要后果的细节，那么规划可能几乎毫无用处。
环境模型应当在其状态和动态两个方面审慎构造，目标是优化规划过程。
模型的各个部分应当持续受到监控，看它们在多大程度上有助于或损害规划效率。
这一领域尚未处理这组复杂问题，也尚未设计出考虑这些影响的模型学习方法。

未来研究需要处理的第四个问题，是自动化地选择智能体所从事的任务，并用这些任务来组织它不断发展的能力。
在机器学习中，通常由人类设计者设定学习智能体预计要掌握的任务。
因为这些任务预先已知并保持固定，所以它们可以被写进学习算法代码中。
但是，展望未来，我们会希望智能体自己选择它应该试图掌握哪些任务。
这些任务可能是某个已经已知的总体任务的子任务；也可能旨在创造构件，使智能体能更高效地学习它未来可能面对、但当前尚未知的许多不同任务。

这些任务可能类似于第 17.1 节讨论的辅助任务或 GVF，也可能是第 17.2 节讨论的由选项解决的任务。
例如，在形成一个 GVF 时，累积信号、策略和终止函数应该是什么？
当前技术水平是人工选择这些东西；但如果能自动做出这些任务选择，就会获得大得多的能力和一般性，尤其是当这些任务来自智能体此前通过表示学习或先前子问题经验所构造出的东西时。
如果 GVF 设计被自动化，那么设计选择本身就必须被显式表示。
任务选择不能再存在于设计者头脑中并写进代码；它们必须以某种方式存在于机器自身中，使它们可以自动设定和改变、监控、过滤并搜索。
任务随后可以像 ANN 中的特征那样层级地构建在其他任务之上。
任务是问题，而 ANN 的内容则是这些问题的答案。
我们预计，将需要一个完整的问题层级，来匹配现代深度学习方法所提供的答案层级。

我们想为未来研究强调的第五个问题，是行为与学习之间通过某种计算上的好奇心类似物而发生的交互。
在本章中，我们一直想象这样一种设定：许多任务使用离策略方法，从同一经验流中同时学习。
采取的动作当然会影响这条经验流，而经验流又会决定发生多少学习以及哪些任务被学到。
当奖励不可得，或并不强烈受行为影响时，智能体可以自由选择在某种意义上使任务学习最大化的动作，也就是说，用某种学习进展的度量作为内部或“内在”奖励，实现一种计算形式的好奇心。
除了衡量学习进展之外，内在奖励还可以表示收到意外的、新奇的或以其他方式有趣的输入，也可以评估智能体在环境中造成变化的能力。
按这些方式生成的内在奖励信号，可以由智能体用来通过定义辅助任务、GVF 或选项而为自己提出任务，如上文所讨论的那样；这样，以这种方式学到的技能就能贡献于智能体掌握未来任务的能力。
结果是一种计算上类似于游戏的东西。
已经有许多关于内在奖励信号这些用途的初步研究；在这一广阔领域中，仍有令人兴奋的未来研究主题。

最后一个需要未来研究关注的问题，是发展方法，使把强化学习智能体嵌入物理环境变得足够安全、可以接受。
这是未来研究最紧迫的领域之一，我们会在下一节进一步讨论它。

\subsection*{17.6\quad 人工智能的未来}

当我们在 20 世纪 90 年代中期撰写本书第一版时，人工智能正在取得显著进展，并正在对社会产生影响；但它大多仍只是激发这些发展的人工智能承诺。
机器学习是那种前景的一部分，但它尚未成为人工智能不可或缺的组成。
到今天，这个承诺已经转化为正在改变数百万人生活的应用，机器学习也已经凭自身成为一项关键技术。
在我们撰写第二版时，人工智能中一些最引人注目的发展涉及强化学习，尤其是“深度强化学习”，即带有深度人工神经网络函数近似的强化学习。
我们正处于人工智能真实世界应用浪潮的开端，其中许多应用会包含强化学习，无论深度与否；它们将以难以预测的方式影响我们的生活。

但是，大量成功的真实世界应用并不意味着真正的人工智能已经到来。
尽管许多领域取得了巨大进展，人工智能与人类智能、甚至与其他动物智能之间的鸿沟仍然很大。
在某些领域，即使像围棋这样可怕的领域，也可以达到超人表现；但发展出像我们一样的系统仍是重大挑战：完整的交互智能体，具有一般适应性和问题求解技能、情感精细性、创造性，以及从经验中快速学习的能力。
强化学习关注通过与动态环境交互来学习；随着未来发展，它将成为具有这些能力的智能体的关键组成部分。

强化学习与心理学和神经科学的联系（第 14 和 15 章）强调了它同人工智能另一个长期目标的相关性：阐明关于心智以及心智如何从大脑中涌现的基本问题。
强化学习理论已经在贡献于我们对大脑奖励、动机和决策过程的理解；也有充分理由相信，通过它与计算精神病学的联系，强化学习理论将对治疗精神障碍的方法作出贡献，包括药物滥用和成瘾。

强化学习未来还能作出的另一种贡献，是帮助人类决策。
在模拟环境中由强化学习得到的策略，可以在人类决策者所处的教育、医疗、交通、能源和公共部门资源分配等领域提供建议。
强化学习的关键特征，即把决策的长期后果纳入考虑，在这里尤其相关。
这一点在双陆棋和围棋等游戏中非常清楚；强化学习的一些最令人印象深刻的结果已经在那里得到展示。
但这也是影响我们生活和地球的许多高风险决策的性质。
强化学习沿袭了过去由许多学科的决策分析师发展出的、为人类决策提供建议的相关方法。
凭借先进的函数近似方法和巨大的计算能力，强化学习方法有潜力克服传统决策支持方法在扩展到更大、更复杂问题时遇到的一些困难。

人工智能快速进步带来了警告：人工智能对我们的社会，甚至对人类自身，构成严重威胁。
著名科学家和人工智能先驱 Herbert Simon 在 2000 年卡内基梅隆大学 Earthware Symposium 的演讲中，预见了我们今天听到的警告。
他谈到任何新知识的承诺与危险之间的永恒冲突，并让我们想起希腊神话中的普罗米修斯和潘多拉：普罗米修斯是现代科学英雄，为了人类利益从诸神那里盗取火；潘多拉的盒子则可能由一个微小而无辜的动作打开，向世界释放无数危险。
Simon 在承认这种冲突不可避免的同时，敦促我们认识到，作为未来的设计者而不是单纯旁观者，我们作出的决定可以让天平向普罗米修斯一边倾斜。
这对强化学习当然也是真的；它能使社会受益，但如果草率部署，也会产生不合意的结果。
因此，涉及强化学习的人工智能应用的\textbf{安全性}是一个值得认真关注的主题。

强化学习智能体可以通过与真实世界交互来学习，也可以通过与真实世界某个片段的模拟交互来学习，或者通过这两种经验来源的混合来学习。
模拟器提供安全环境，使智能体可以探索和学习，而不冒着对自身或环境造成真实损害的风险。
在当前大多数应用中，策略从模拟经验中学习，而不是从同真实世界的直接交互中学习。
除了避免不合意的真实世界后果，从模拟经验中学习还可以为学习提供几乎无限的数据；通常成本低于获得真实经验所需成本；而且因为模拟通常比真实时间快得多，学习往往也能比依赖真实经验更快发生。

然而，强化学习的全部潜力要求强化学习智能体嵌入真实世界经验流中；它们在我们的世界中行动、探索和学习，而不只是它们自己的世界中。
毕竟，强化学习算法，至少是本书关注的那些算法，被设计成线上学习；它们在许多方面模仿动物如何在非平稳且敌对的环境中生存。
把强化学习智能体嵌入真实世界，在实现人工智能增强和扩展人类能力的承诺方面，可能具有变革性。

希望强化学习智能体在真实世界中行动和学习的一个主要原因是：要以足够高的保真度模拟真实世界经验，使所得策略，无论由强化学习还是其他方法得到，在指挥真实动作时表现良好并且安全，常常很困难，有时甚至不可能。
这对那些动态依赖人类行为的环境尤其真实，例如教育、医疗、交通和公共政策；这些领域肯定能从改进的决策中受益。
然而，关于人工智能潜在危险的警告最需要被嵌入真实世界的智能体留意。

这些警告中的一些与强化学习特别相关。
因为强化学习基于优化，它继承了所有优化方法的优点和缺点。
缺点方面的问题是设计目标函数，或者在强化学习中设计奖励信号，使优化产生期望结果并避免不合意结果。
我们在第 17.4 节说过，强化学习智能体可能发现意外方式让环境提供奖励；其中一些方式可能不合意，甚至危险。
当我们只间接地指定希望系统学会什么时，正如设计强化学习系统的奖励信号时那样，在学习完成之前，我们不会知道智能体会多接近地满足我们的愿望。
这对强化学习来说并不是什么新问题；在文学和工程中，人们很早就认识到了它。
例如，在歌德的诗《魔法师的学徒》中，学徒用魔法让扫帚为他取水，但由于学徒对魔法知识不足，结果造成了意外洪水。
在工程语境中，控制论创始人 Norbert Wiener 半个多世纪前就通过引用“猴爪”的超自然故事警告过这个问题：“……它给你的，是你要求的东西，而不是你本该要求的东西，也不是你意图的东西”（第 59 页）。
Nick Bostrom（2014）也在现代语境中详细讨论过这个问题。
任何有强化学习经验的人，都可能见过系统发现意外方式来获得大量奖励。
有时，意外行为是好的：它以漂亮的新方式解决了问题。
在另一些情形中，智能体学到的东西违反了系统设计者也许从未想到过的考虑。
如果一个智能体要在真实世界行动，并且没有机会让人类审查其行动或简单中断其行为，那么仔细设计奖励信号是必不可少的。

尽管可能有意外的负面后果，优化已经被工程师、建筑师以及其他设计对世界产生积极影响的人使用了数百年。
我们的环境中有许多好东西都归功于优化方法的应用。
已经发展出许多方法来降低优化风险，例如加入硬约束和软约束，把优化限制到稳健和风险敏感的策略，以及使用多个目标函数进行优化。
其中一些方法已经被改造用于强化学习；为了处理这些关切，还需要更多研究。
确保强化学习智能体的目标与我们的目标协调一致，仍是一个挑战。

如果强化学习智能体要在真实世界中行动和学习，另一个挑战不仅在于它们最终可能学到什么，也在于它们学习时会如何行动。
怎样确保智能体获得足够经验以学到高表现策略，同时又不伤害其环境、其他智能体或自身（或者更现实地说，把伤害概率保持在可接受的低水平）？
这个问题也并非强化学习所独有，也不是新问题。
嵌入式强化学习的风险管理与缓解，类似于控制工程师从一开始在使用自动控制时就不得不面对的问题：控制器行为可能产生不可接受的、甚至灾难性的后果，例如控制飞机或精细化工过程时。
控制应用依赖仔细的系统建模、模型验证和广泛测试；也有高度发展的理论，旨在当被控系统的动态并非完全已知时，保证自适应控制器的收敛和稳定性。
理论保证从不坚不可摧，因为它们依赖于数学背后假设的有效性；但如果没有这种理论，以及风险管理和缓解实践，自动控制，无论自适应与否，就不会像今天这样有益于改善我们所依赖过程的质量、效率和成本效益。
未来强化学习研究最紧迫的领域之一，是适应并扩展控制工程中发展出来的方法，目标是让强化学习智能体能够以可接受的安全性完全嵌入物理环境。

最后，我们回到 Simon 对我们的呼吁：认识到我们是未来的设计者，而不是单纯旁观者。
通过我们作为个体所作的决定，以及我们能够对社会治理方式施加的影响，我们可以努力确保新技术可能带来的益处超过它可能造成的伤害。
在强化学习的情形中，有充分机会做到这一点；强化学习可以帮助改善地球上生活的质量、公平性和可持续性，但也可能释放新的危险。
一个已经到来的威胁是人工智能应用造成的工作岗位替代。
尽管如此，仍有充分理由相信，人工智能的好处可以超过它造成的扰动。
至于安全性，强化学习可能带来的危险并不完全不同于那些已经在优化和控制方法的相关应用中被成功管理的危险。
随着强化学习在未来应用中走向真实世界，开发者有义务遵循类似技术中已经演化出的最佳实践，同时扩展这些实践，以确保普罗米修斯占上风。

\section*{书目与历史说明}

\noindent\textbf{17.1}\quad 一般价值函数首先由 Sutton 及其同事明确识别出来。
Ring 发展了一套关于 GVF（“预报”）的、虽尚未发表却颇有影响的广泛实验思想。

强化学习中多头学习的最早展示来自 Jaderberg 等人（2017）。
Bellemare、Dabney 和 Munos（2017）表明，预测奖励分布中更多内容可以显著加快对其期望的优化学习；这是辅助任务的一个例子。
此后许多人沿着这条研究路线继续推进。

据我们所知，把经典条件化的一般理论理解为学到的预测加上对这些预测的内置反射反应，在心理学文献中尚未被清楚阐明。
Modayil 和 Sutton（2014）把它描述为一种机器人和其他智能体的工程方法，并称其为“巴甫洛夫控制”，以暗示它在经典条件化中的根源。

\noindent\textbf{17.2}\quad 把时间抽象的行动过程形式化为选项，是 Sutton、Precup 和 Singh（1999）引入的；它建立在 Parr（1998）和 Sutton（1995a）的先前工作以及关于半 MDP 的经典工作之上。
Precup（2000）的博士论文充分发展了选项思想。
这些早期工作的一个重要限制是，它们没有处理带函数近似的离策略情形。
选项内学习通常需要离策略学习，而当时用函数近似还不能可靠做到这一点。
虽然现在已有多种使用函数近似的稳定离策略学习方法，但在本书出版时，它们与选项思想的结合尚未得到显著探索。
Barto 和 Mahadevan（2003）以及 Hengst（2012）综述了选项形式化和其他时间抽象方法。

用 GVF 实现选项模型此前没有被描述过。
我们的表述使用了 Modayil、White 和 Sutton（2014）为预测策略终止处信号而引入的技巧。

少数学习带函数近似的选项模型的工作，包括 Sorg 和 Singh（2010）以及 Bacon、Harb 和 Precup（2017）的工作。

选项和选项模型向平均奖励设定的扩展，在文献中尚未发展出来。

\noindent\textbf{17.3}\quad Monahan（1982）给出了 POMDP 方法的良好介绍。
PSR 和测试由 Littman、Sutton 和 Singh（2002）引入。
OOM 由 Jaeger（1997，1998，2000）引入。
序贯系统统一了 PSR、OOM 和许多其他工作，由 Michael Thon 的博士论文引入（2017；Thon 和 Jaeger，2015）。
时间关系网络的扩展由 Tanner（2006；Sutton 和 Tanner，2005）发展，随后扩展到选项（Sutton、Rafols 和 Koop，2006）。

带非马尔可夫状态表示的强化学习理论，由 Singh、Jaakkola 和 Jordan（1994；Jaakkola、Singh 和 Jordan，1995）明确发展。
早期面向部分可观测性的强化学习方法，由 Chrisman（1992）、McCallum（1993，1995）、Parr 和 Russell（1995）、Littman、Cassandra 和 Kaelbling（1995），以及 Lin 和 Mitchell（1992）发展。

\noindent\textbf{17.4}\quad 在强化学习中包含建议和教学的早期尝试，包括 Lin（1992）、Maclin 和 Shavlik（1994）、Clouse（1996）以及 Clouse 和 Utgoff（1992）的工作。

Skinner 的塑形不应与 Ng、Harada 和 Russell（1999）引入的“基于势的塑形”技术混淆。
Wiewiora（2003）已经表明，他们的技术等价于提供价值函数的初始近似这一更简单思想，如（17.10）所示。

\noindent\textbf{17.5}\quad 关于当今深度学习技术的讨论，我们推荐 Goodfellow、Bengio 和 Courville（2016）的书。
ANN 中灾难性干扰问题由 McCloskey 和 Cohen（1989）、Ratcliff（1990）以及 French（1999）发展。
回放缓冲区思想由 Lin（1992）引入，并在 Atari 游戏系统的深度学习中被突出使用（第 16.5 节；Mnih 等，2013，2015）。

Minsky（1961）是最早识别表示学习问题的人之一。

少数考虑用学到的近似模型进行规划的工作，包括 Kuvayev 和 Sutton（1996）、Sutton、Szepesvari、Geramifard 和 Bowling（2008）、Nouri 和 Littman（2009），以及 Hester 和 Stone（2012）的工作。

为了避免拖慢规划，在模型构造中需要具有选择性，这在人工智能中是众所周知的。
一些经典工作来自 Minton（1990）以及 Tambe、Newell 和 Rosenbloom（1990）。
Hauskrecht、Meuleau、Kaelbling、Dean 和 Boutilier（1998）在带确定性选项的 MDP 中展示了这种效应。

Schmidhuber（1991a,b）提出，如果奖励信号是智能体环境模型改进速度的函数，那么就会产生类似好奇心的东西。
Klyubin、Polani 和 Nehaniv（2005）提出的赋权函数，是对智能体控制其环境能力的信息论度量，可以作为内在奖励信号。
Baldassarre 和 Mirolli（2013）收录了研究内在奖励和动机的学者从生物与计算两个角度作出的贡献，其中包括关于“内在动机强化学习”的观点；这一术语由 Singh、Barto 和 Chentenez（2004）引入。
另见 Oudeyer 和 Kaplan（2007）、Oudeyer、Kaplan 和 Hafner（2007）以及 Barto（2013）。

\end{document}
